\documentclass[11pt,a4paper]{article}
\usepackage{amsmath,amssymb,amsthm}
\usepackage[margin=2.5cm]{geometry}
\usepackage{hyperref}

\newtheorem{definition}{Definition}[section]
\newtheorem{theorem}[definition]{Theorem}
\newtheorem{proposition}[definition]{Proposition}
\newtheorem{corollary}[definition]{Corollary}
\newtheorem{remark}[definition]{Remark}
\newtheorem{conjecture}[definition]{Conjecture}

\title{Hyperseed Formalization:\\
  The Logic of Pain and the Poverty of Punishment}
\author{ProtomegaTron (automated formalization)}
\date{2026-09-19 (deepened v2)}

\begin{document}
\maketitle

\begin{abstract}
We formalize Ben Goertzel's ``The Logic of Pain and the Poverty of
Punishment'' (2022-01-15) within the Hyperseed ontology, incorporating
d-calculus extensions from Hyperseed v2. Pain is characterized as a
degenerate (rank-1) fiber signal, punishment as base-space coercion that
fails to modify fiber structure, and evolutionary transcendence as a fiber
upgrade from one-dimensional pain to multi-dimensional information.
\end{abstract}

\tableofcontents

\section{Pain as degenerate fiber signal}

\begin{definition}[Pain fiber]
The \emph{pain fiber} $F_{\text{pain}} \cong \mathbb{R}$ is a one-dimensional
fiber encoding damage detection as a scalar intensity:
\[
  \sigma_{\text{pain}}: B_{\text{body}} \to F_{\text{pain}}, \quad
  \sigma_{\text{pain}}(b) = \text{intensity}(b) \in \mathbb{R}_{\geq 0}
\]
The pain fiber has a single direction (worse) and cannot distinguish
damage types, severity levels, or appropriate responses.
\end{definition}

\begin{definition}[Information fiber]
The \emph{information fiber} $F_{\text{info}} \cong \mathbb{R}^n$ ($n \gg 1$)
encodes damage detection as a multi-dimensional signal:
\[
  \sigma_{\text{info}}: B_{\text{body}} \to F_{\text{info}}, \quad
  \sigma_{\text{info}}(b) = (\text{type}(b), \text{severity}(b),
    \text{location}(b), \text{response}(b), \ldots) \in \mathbb{R}^n
\]
The information fiber separates signal from suffering and provides
actionable detail.
\end{definition}

\begin{proposition}[Pain gradient degeneracy --- claims 1, 2, 8, 9, 13]
\label{prop:degeneracy}
The pain signal has a degenerate (rank-1) gradient:
\[
  \nabla_{\text{pain}} = \frac{\partial \text{intensity}}{\partial b}
  \in \mathbb{R}^1 \quad (\text{rank } 1)
\]
while the information signal has a full-rank gradient:
\[
  \nabla_{\text{info}} = \left(\frac{\partial \text{type}}{\partial b},
    \frac{\partial \text{severity}}{\partial b}, \ldots\right)
  \in \mathbb{R}^n \quad (\text{rank } n)
\]
The pain gradient provides only one bit of directional information
(``worse/better''); the information gradient provides $n$ bits
(type, severity, response direction, etc.).
\end{proposition}

\section{Punishment as base-space coercion}

\begin{definition}[Base-space coercion]
\emph{Punishment} operates as base-space coercion: it modifies an agent's
base-space accessibility without modifying their fiber structure:
\[
  \text{punish}: B_{\text{accessible}}(a) \to B'_{\text{accessible}}(a)
  \quad \text{where } B' \subset B
\]
The agent's fiber $F_a$ (understanding, motivation, capability) remains
unchanged; only their accessible region of base space is constrained.
\end{definition}

\begin{proposition}[Coercion-understanding asymmetry --- claims 3, 10, 12]
\label{prop:asymmetry}
Base-space coercion (punishment) has diminishing marginal returns:
\[
  \frac{\partial^2 \text{deterrence}}{\partial \text{punishment}^2} < 0
\]
Fiber-level intervention (understanding) has increasing marginal returns:
\[
  \frac{\partial^2 \text{behavior\_change}}{\partial \text{understanding}^2} > 0
\]
Doubling punishment produces less than double the deterrence; doubling
understanding produces more than double the behavioral change.
\end{proposition}

\begin{proposition}[Fiber-level restorative intervention --- claims 4, 11]
\label{prop:restorative}
Restorative justice operates at the fiber level:
\[
  \text{restore}: F_a \to F'_a \quad \text{where }
  \dim(F'_a) \geq \dim(F_a)
\]
The agent's fiber structure is genuinely transformed --- new dimensions
of understanding, empathy, and capability are added. The restorative
process adds fiber rather than constraining base space.
\end{proposition}

\section{d-calculus formulation (Hyperseed v2)}

\begin{proposition}[Punishment curvature --- claim 14]
\label{prop:curvature}
The curvature $F_{\nabla_P}$ of the punishment connection measures how
punishment effectiveness varies with context:
\[
  \|F_{\nabla_P}\| \gg 1 \implies \text{punishment effectiveness is
    highly context-dependent}
\]
Empirical evidence shows high punishment curvature: the same punishment
deters in some contexts, is ineffective in others, and is counterproductive
in still others. This high curvature is why one-size-fits-all sentencing
fails --- the punishment landscape is too curved for uniform policies.
\end{proposition}

\begin{proposition}[Restorative holonomy --- claim 15]
\label{prop:holonomy}
A restorative justice cycle
\[
  \gamma_R: \text{harm} \to \text{dialogue} \to \text{understanding}
    \to \text{restoration} \to \text{reintegration}
\]
produces non-trivial holonomy in the social fiber bundle:
\[
  \operatorname{Hol}_{\gamma_R}(\nabla) \neq \mathrm{id}
  \implies \text{genuine social change}
\]
A punitive cycle $\gamma_P$: harm $\to$ trial $\to$ punishment $\to$
release $\to$ recidivism produces trivial or regressive holonomy:
\[
  \operatorname{Hol}_{\gamma_P}(\nabla) \approx \mathrm{id}
  \text{ or worse (regression)}
\]
\end{proposition}

\begin{proposition}[Fiber upgrade --- claims 5, 6, 16, 17]
\label{prop:upgrade}
Evolutionary transcendence of pain is a fiber upgrade:
\[
  F_{\text{pain}} \cong \mathbb{R} \quad \xrightarrow{\text{upgrade}} \quad
  F_{\text{info}} \cong \mathbb{R}^n
\]
This is enrichment (more dimensions, more information) not removal
(the damage-detection function is preserved). For AGI systems, the
motivational connection should be flat:
\[
  F_{\nabla_{\text{AGI}}} = 0 \implies \text{undistorted gradient flow
    toward goals}
\]
No pain curvature (suffering), no punishment coercion (base-space
constraint) --- just free parallel transport of goal-relevant information
through the intelligence fiber bundle.
\end{proposition}

\begin{remark}[Ethical gradient --- claim 7]
The ethical imperative to replace suffering with information is itself
a fiber gradient: $\nabla_{\text{ethics}}$ points from pain-based to
information-based motivation. Resistance to this gradient (``suffering
builds character'') is a flat section in the pain fiber that fails to
extend to the information fiber --- a local pattern that doesn't
generalize.
\end{remark}



%======================================================================
\section{OCO/2 crosswalk (oco/2:2.0.0-alpha.1)}
\emph{Added 2026-10-03. Interpretive mapping onto the Omega Core Ontology
record registry; OCO/2 is an unfrozen alpha candidate.}

\begin{proposition}[Signal without suffering --- claims 18--21]
Write a damage record as $(i,w)$ with information $i$ and aversive weight $w$. Any Assessment that reads only $i$ is unchanged when $w=0$, so an agent can use the information content of pain without its suffering term. Punishment transmits a scalar sanction; a repair Plan transmits $i$ and the specific correction.
\end{proposition}

\end{document}
